\subsection{Independence}
\hrulefill

\paragraph*{Definition:}
Two events $A$ and $B$ are independent if the occurance of one event does not affect the probability of the other event occurring.
Otherwise, the events are dependent. Sometimes it is obvios that two events influence each other, such as crude oil prices and 
gas prices. Sometimes it is obvious that two events are not dependent, such as rolling a die and the weather.
Sometimes it is difficult to tell, such as obtaining two heads in three coin flips and obtaining a head on the first of three flips.
Mutual exclusivity and independence are not the same thing. Two events can be independent and not mutually exclusive, such as rolling 
a 1 and rolling an odd number on a die.

\paragraph*{Checking for Independence:}
In many cases, you will need to evaluate independence by checking the following formulas:
\begin{itemize}
    \item A and B are independent if and only if $P(A|B) = P(A)$
    \item A and B are independent if and only if $P(A \cap B) = P(A)P(B)$
\end{itemize}

Note that independent trials are okay to be assumed to be independent, such as rolling a die or flipping a coin.

\paragraph*{Properties:}
\begin{itemize}
    \item If A and B are independent, then $b$ and $A$ are independent.
    \item If A and B are independent, then $A'$ and $B'$, $A$ and $B'$, and $A'$ and $B$ are all independent.
\end{itemize}

\pagebreak

\paragraph*{Example:}
Suppose we draw a card from a deck of 52 playing cards. Let A be the event that the card is an Ace, and let B be the event that 
the card is a Spade. Are A and B independent?

\begin{align*}
    P(A) &= \frac{4}{52} = \frac{1}{13} \\
    P(B) &= \frac{13}{52} = \frac{1}{4} \\
    P(A \cap B) &= \frac{1}{52} \\
    P(A)P(B) &= \frac{1}{13} \cdot \frac{1}{4} = \frac{1}{52} \\
\end{align*}
Since $P(A \cap B) = P(A)P(B)$, A and B are independent.

\paragraph*{Example:}
Suppose we roll 2d6. What is the probability of getting 2 even numbers?
\begin{align*}
    P(E_1  \cap E_2) &= P(E_1)P(E_2) &\text{ since } E_1 \text{ and } E_2 \text{ are independent} \\
    &= \left(\frac{3}{6}\right)\left(\frac{3}{6}\right) \\
    &= \frac{9}{36} = \frac{1}{4}
\end{align*}

\paragraph*{Example:}
Roll 2d6. Let A be the event that the sum is 7. Let B be the event that the first die is a 4. Are A and B independent?
\begin{align*}
    P(A) &= \frac{6}{36} = \frac{1}{6} \\
    P(B) &= \frac{1}{6} \\
    P(A \cap B) &= \frac{1}{36} \\
    P(A)P(B) &= \frac{1}{6} \cdot \frac{1}{6} = \frac{1}{36} \\
\end{align*}
Since $P(A \cap B) = P(A)P(B)$, A and B are independent.